\section{Methods}

\subsection{FEMA NFHL Corridor Overlay}

For each of the 227 interstate corridors in the ROUTE v1.2 corpus, we extract the corridor
centerline geometry from the TIGER/Line 2023 dataset \citep{ROUTE_A1}. We construct a 1-mile
bilateral buffer (2-mile corridor envelope) using the GEOS library (PostGIS implementation),
then intersect the buffered corridor against FEMA NFHL SFHA polygons retrieved via the FEMA
Map Service Center REST API. The intersection produces, for each corridor:

\begin{itemize}
\item \textbf{Total SFHA miles}: sum of all buffered corridor miles that overlap any SFHA
polygon, regardless of contiguity.
\item \textbf{Maximum consecutive SFHA miles}: length of the longest contiguous run of
SFHA exposure within the corridor buffer, measured along the corridor centerline after
projecting SFHA polygon membership to 0.1-mile corridor segments.
\end{itemize}

The consecutive-miles metric is computed as the longest sequence of consecutive 0.1-mile
segments for which the corridor buffer intersects an SFHA polygon. This operationalizes
structural exposure: a corridor with 127 consecutive SFHA miles is one in which 127 miles
of road, drainage infrastructure, and access roads are all within the 100-year flood plain
simultaneously --- meaning the entire sub-corridor floods as a unit during a 100-year event.

\subsection{D1 Composite Score Formula}

The D1 score for flood-exposed corridors combines consecutive and total SFHA miles under a
weighted formula calibrated to the ROUTE rubric version 1.2 \citep{ROUTE_D1}:

\begin{equation}
D1_{\text{flood}} = 0.7 \times f_c(m_c) + 0.3 \times f_t(m_t)
\end{equation}

where $m_c$ is maximum consecutive SFHA miles, $m_t$ is total SFHA miles, and $f_c$, $f_t$
are piecewise-linear scoring functions anchored as follows:

\begin{align}
f_c(m_c) &= \begin{cases}
0 & m_c = 0 \\
5 & m_c = 15 \text{ mi} \\
10 & m_c \geq 60 \text{ mi}
\end{cases} \\[6pt]
f_t(m_t) &= \begin{cases}
0 & m_t = 0 \\
5 & m_t = 40 \text{ mi} \\
10 & m_t \geq 150 \text{ mi}
\end{cases}
\end{align}

Values between anchor points are linearly interpolated. Values at or above the upper anchor
receive a score of 10 (capped). The 70\%/30\% weighting reflects the greater investment
significance of consecutive exposure: a corridor where 60+ consecutive miles are in the
flood plain faces a qualitatively different hardening challenge than a corridor where the
same 60 miles are fragmented across non-contiguous flood-zone segments.

\subsection{Winter Closure Scoring}

For corridors where the primary climate exposure mechanism is winter precipitation (snowpack,
ice, avalanche) rather than flood zone location, we compute D1 using closure frequency data:

\begin{equation}
D1_{\text{winter}} = f_w\!\left(\text{closures/yr} \times \text{mean duration (hr)}\right)
\end{equation}

where $f_w$ maps annual closure-hours to a 0--10 score anchored on: 0 closure-hours $\to$ 0;
450 closure-hours $\to$ 5; 900 closure-hours $\to$ 10. Donner Pass at 50 closures/yr $\times$
18 hr mean = 900 closure-hours/yr maps to D1 = 7.8 after the calibration adjustment for the
fact that winter closures are recoverable (hours to days) while flood closures may persist
through a multi-day event.

For corridors with dual exposure (e.g., I-5 Siskiyou with both wildfire smoke and winter
precipitation), we score each mechanism separately and take the maximum, reflecting that the
harder-to-mitigate mechanism drives investment needs.

\subsection{2050 Projection Methodology}
\label{sec:projections}

We generate 2050 D1 projections for coastal corridors by applying NOAA Intermediate SLR
(0.5m local-adjusted for Gulf Coast) to the 1/3 arc-second National Elevation Dataset. The
procedure:

\begin{enumerate}
\item Extract elevation raster for a 5-mile corridor envelope.
\item Apply projected SLR offset to the current-day SFHA inundation threshold (1\% annual
exceedance elevation at each NFHL cross-section).
\item Recompute SFHA polygon membership for each 0.1-mile corridor segment at the projected
flood level.
\item Recalculate $m_c$ and $m_t$ and apply the D1 formula to produce the 2050 projected score.
\end{enumerate}

Land subsidence is incorporated for Louisiana corridors (1.5 cm/yr median, equivalent to
+0.375m additional effective SLR by 2050 relative to the NOAA global figure) and Houston
subsidence zone (0.8 cm/yr) using USGS InSAR-derived subsidence maps. This produces Gulf
Coast effective local SLR of 0.875m by 2050 under the NOAA Intermediate global scenario.
